% 87-02 ④(87쪽) — 보기 그래프: $x=-1$ 에서 기울기 0, $x=3$ 에서 극소인 아래로 볼록한 사차함수 개형. 그림 자체가 보기라 TikZ 로 다시 그림.
% 라벨은 원문 그대로 -1, O, 3, x, y, y=f(x) 와 점선 두 줄만.
\begin{tikzpicture}[baseline=(current bounding box.center), x=0.36cm, y=0.36cm, line width=0.5pt]
  \path (-2.9,-2.4) rectangle (5.15,3.5);
  \draw[->, >=stealth, line width=0.7pt] (-2.7,0) -- (4.85,0) node[below=1pt, font=\footnotesize] {$x$};
  \draw[->, >=stealth, line width=0.7pt] (0,-2.0) -- (0,3.2) node[left=1pt, font=\footnotesize] {$y$};
  \draw[densely dotted, line width=0.4pt] (-1,0) -- (-1,0.55);
  \draw[densely dotted, line width=0.4pt] (3,0) -- (3,-1.02);
  \draw[line width=0.6pt, smooth] plot coordinates {(-2.4,2.9) (-2.25,2.35) (-2.1,1.85) (-1.9,1.35) (-1.7,1.0) (-1.5,0.78) (-1.3,0.63) (-1,0.55) (-0.7,0.45) (-0.4,0.3) (-0.1,0.12) (0.3,-0.15) (0.7,-0.42) (1.2,-0.7) (1.7,-0.88) (2.2,-0.98) (2.7,-1.02) (3,-1.02) (3.3,-0.95) (3.6,-0.7) (3.85,-0.35) (4.05,0.1) (4.25,0.8) (4.4,1.5) (4.5,2.15)};
  \node[below=1pt, font=\footnotesize] at (-1.4,0) {$-1$};
  \node[below left=1pt and 1pt, font=\footnotesize] at (0,0) {$\pt{O}$};
  \node[above=1pt, font=\footnotesize] at (3,0) {$3$};
  \node[anchor=west, font=\footnotesize] at (0.25,2.95) {$y=f(x)$};
\end{tikzpicture}
